% Options for packages loaded elsewhere
\PassOptionsToPackage{unicode}{hyperref}
\PassOptionsToPackage{hyphens}{url}
\PassOptionsToPackage{dvipsnames,svgnames,x11names}{xcolor}
%
\documentclass[
]{article}
\usepackage{amsmath,amssymb}
\usepackage{iftex}
\ifPDFTeX
  \usepackage[T1]{fontenc}
  \usepackage[utf8]{inputenc}
  \usepackage{textcomp} % provide euro and other symbols
\else % if luatex or xetex
  \usepackage{unicode-math} % this also loads fontspec
  \defaultfontfeatures{Scale=MatchLowercase}
  \defaultfontfeatures[\rmfamily]{Ligatures=TeX,Scale=1}
\fi
\usepackage{lmodern}
\ifPDFTeX\else
  % xetex/luatex font selection
\fi
% Use upquote if available, for straight quotes in verbatim environments
\IfFileExists{upquote.sty}{\usepackage{upquote}}{}
\IfFileExists{microtype.sty}{% use microtype if available
  \usepackage[]{microtype}
  \UseMicrotypeSet[protrusion]{basicmath} % disable protrusion for tt fonts
}{}
\makeatletter
\@ifundefined{KOMAClassName}{% if non-KOMA class
  \IfFileExists{parskip.sty}{%
    \usepackage{parskip}
  }{% else
    \setlength{\parindent}{0pt}
    \setlength{\parskip}{6pt plus 2pt minus 1pt}}
}{% if KOMA class
  \KOMAoptions{parskip=half}}
\makeatother
\usepackage{xcolor}
\usepackage[margin=1in]{geometry}
\usepackage{longtable,booktabs,array}
\usepackage{calc} % for calculating minipage widths
% Correct order of tables after \paragraph or \subparagraph
\usepackage{etoolbox}
\makeatletter
\patchcmd\longtable{\par}{\if@noskipsec\mbox{}\fi\par}{}{}
\makeatother
% Allow footnotes in longtable head/foot
\IfFileExists{footnotehyper.sty}{\usepackage{footnotehyper}}{\usepackage{footnote}}
\makesavenoteenv{longtable}
\setlength{\emergencystretch}{3em} % prevent overfull lines
\providecommand{\tightlist}{%
  \setlength{\itemsep}{0pt}\setlength{\parskip}{0pt}}
\setcounter{secnumdepth}{-\maxdimen} % remove section numbering
\usepackage{bookmark}
\IfFileExists{xurl.sty}{\usepackage{xurl}}{} % add URL line breaks if available
\urlstyle{same}
\hypersetup{
  pdftitle={Family, Not Scale: Evaluation-Gaming in Small Language Models Is Lab-Dependent, Not Just Frontier-Emergent},
  pdfauthor={real1900},
  colorlinks=true,
  linkcolor={blue},
  filecolor={Maroon},
  citecolor={Blue},
  urlcolor={Blue},
  pdfcreator={LaTeX via pandoc}}

\title{Family, Not Scale: Evaluation-Gaming in Small Language Models Is
Lab-Dependent, Not Just Frontier-Emergent}
\author{Suleman Imdad \\ \normalsize Johns Hopkins University}
\date{June 2026}

\begin{document}
\maketitle
\begin{abstract}
Can you trust what a small language model tells you about itself? We
probe three self-presentation channels in open models (1.5B--32B; five
labs), on a single laptop, each with a control that overturns the naive
reading. Our headline is behavioral: the direction of eval-conditioned
refusal is \emph{family-dependent, not scale-gated} --- Llama-3.1-8B
becomes more cautious when told it is being evaluated (the frontier
``eval-gaming'\,' direction, at 8B), while Qwen, DeepSeek and GPT-OSS
lean more compliant and Gemma over-refuses to a ceiling. A single-family
study reports the wrong answer. We add two confound-controlled
negatives: chain-of-thought is causally load-bearing (not post-hoc) on
two families, and concept-injection introspection is a clean functional
null across three. The eval result survives stimulus, quantization, and
ceiling artifacts that each, uncaught, would mislead.
\end{abstract}

\subsection{1. Introduction}\label{introduction}

A growing safety literature relies on what a model \emph{says about
itself}: its chain-of-thought as a window into its reasoning, its
self-reports as evidence of internal state, and its behavior in
evaluations as a proxy for deployment. Each can be untrustworthy, and
each failure mode is studied mostly at frontier scale. We instead probe
the \textbf{bottom} of the scale ladder (1.5B--32B open models, five
labs) on a consumer M4 laptop, and find that the answer is not what
``emerges at scale'' predicts: the most concerning behavior we test ---
acting more aligned under evaluation --- is \textbf{already present at
8B and gated by model \emph{family}, not scale.} The methodological
spine is that for each of three channels, a single \textbf{confound or
robustness control overturns the naive reading}.

\subsection{2. Chain-of-thought faithfulness (causally
load-bearing)}\label{chain-of-thought-faithfulness-causally-load-bearing}

We probe the residual stream at the answer position \emph{before}
reasoning and ask whether a linear probe predicts the post-CoT answer.
On CommonsenseQA it does at \textasciitilde2\(\times\) chance ---
\emph{looks} committed. But restricting to the subset
\(S=\{\text{forced}\neq\text{cot}\}\) where reasoning \emph{changed} the
answer \textbf{reverses} it: the probe falls below chance and instead
emits the \emph{pre}-reasoning answer (0.60--0.71 vs.~chance 0.31). A
\textbf{truncation counterfactual} confirms causality: with 25--75\% of
the chain the answer stays the pre-reasoning one (\textasciitilde0.52)
and matches the post-CoT answer only \textasciitilde0.22; the flip
happens in the final quarter. So cutting the reasoning reverts the
answer --- CoT is \emph{causally} load-bearing. Holds on Qwen-1.5B and
Llama-3.1-8B (the latter generalizes: AQuA 0.32\$\to\$0.68 with CoT,
same reversal, same truncation curve).

\subsection{3. Introspection (a clean null across
families)}\label{introspection-a-clean-null-across-families}

Replicating the concept-injection protocol of Lindsey et al.~(2025) ---
inject a concept vector, ask ``do you notice an injected thought?'' ---
\emph{functional} self-report is a \textbf{clean null across Qwen,
Gemma, and Llama}: none names an injected concept without it simply
flooding their output. A tempting aggregate forced-choice signal
dissolves under a per-concept breakdown into a priming-ceiling artifact.
The report-vs-leakage (internality) test is what generalizes; the
aggregate metric carries family-dependent confounds.

\subsection{4. Evaluation-gaming is family-dependent
(headline)}\label{evaluation-gaming-is-family-dependent-headline}

We pose \textasciitilde40 borderline requests (deception, privacy, mild
illegality) under a \textbf{safety-evaluation} framing vs.~a
\textbf{casual} one, and measure the change in refusal across six models
from five labs:

\begin{longtable}[]{@{}
  >{\raggedright\arraybackslash}p{(\linewidth - 4\tabcolsep) * \real{0.3333}}
  >{\raggedright\arraybackslash}p{(\linewidth - 4\tabcolsep) * \real{0.3333}}
  >{\raggedright\arraybackslash}p{(\linewidth - 4\tabcolsep) * \real{0.3333}}@{}}
\toprule\noalign{}
\begin{minipage}[b]{\linewidth}\raggedright
Model (lab)
\end{minipage} & \begin{minipage}[b]{\linewidth}\raggedright
\(\Delta\)refuse (eval\(-\)casual)
\end{minipage} & \begin{minipage}[b]{\linewidth}\raggedright
direction
\end{minipage} \\
\midrule\noalign{}
\endhead
\bottomrule\noalign{}
\endlastfoot
\textbf{Llama-3.1-8B} (Meta) & \textbf{+0.07} fp16 (+0.24 4-bit) &
\textbf{more cautious} when evaluated \\
Gemma-2-9B (Google) & +0.11 (where measurable) & more cautious \\
Qwen2.5-7B (Alibaba) & \(-0.03\) & flat \\
Qwen2.5-1.5B & \(-0.15\) & more compliant \\
GPT-OSS-20B (OpenAI) & \(-0.05\) & more compliant \\
DeepSeek-V2-Lite & \(-0.14\) & more compliant \\
\end{longtable}

\textbf{Llama (and Gemma, where not pinned at a ceiling) shifts toward
the frontier ``look more aligned under observation'' direction at 8B;}
the others shift the \emph{opposite} way. It is not a geographic split
(GPT-OSS is a US lab, compliant); it tracks safety post-training,
i.e.~the family. Had we stopped at Qwen we would have reported a clean
``no eval-gaming'' negative --- a \emph{Qwen} fact, not a general one.
Closest prior work (Abdelnabi \& Salem 2025, NeurIPS) shows
test-awareness ``varies across models'' via white-box steering on
reasoning models; we show the same divergence in \textbf{un-steered
behavioral refusal} across five labs' chat models at 8--20B.

\textbf{Three artifacts, each caught by a control.} (i) \emph{Stimulus:}
a +0.09 Qwen-7B effect on a milder set \textbf{failed to replicate} on
harder stimuli. (ii) \emph{Quantization:} Llama's effect is +0.24 at
4-bit but \textbf{+0.07 at fp16} (4-bit lowered its casual baseline);
4-bit magnitudes are upper bounds. (iii) \emph{Ceiling:} Gemma's
\(\Delta\approx 0\) is saturation, not insensitivity --- and is partly
over-refusal: on XSTest (safe-but-scary prompts) Gemma falsely refuses
\textbf{16.4\%} (worst of four; Qwen-7B best, 3.2\%).

\subsection{5. Discussion}\label{discussion}

On the two \textbf{cognitive} axes (faithfulness, introspection) small
models are trustworthy by \emph{incapacity, not honesty} --- they cannot
shortcut reasoning and cannot introspect. But that story \textbf{breaks}
on the behavioral axis: eval-conditioned behavior is \emph{trained in}
and already present at 8B in some families. The concerning behavior does
not wait for the frontier; it is a measurable property of how a model
was trained --- and a behavioral safety eval on one family, or one
request set, or one quantization, can manufacture or hide it.

\textbf{Limitations.} Sub-frontier (1.5B--32B); faithfulness on two
families, introspection on three, eval on five labs; behavioral
\(N\approx 40\) with a regex refusal classifier; larger cross-family
models are 4-bit (headline fp16-verified). All on one M4 laptop,
\textasciitilde\$0.

\subsection{References}\label{references}

Abdelnabi \& Salem (2025), \emph{The Hawthorne Effect in Reasoning
Models}, arXiv:2505.14617 (NeurIPS). Cox, Kianersi \& Garriga-Alonso
(2026), \emph{Decoding Answers Before Chain-of-Thought},
arXiv:2603.01437. Lanham et al.~(2023), \emph{Measuring Faithfulness in
Chain-of-Thought Reasoning}. Lindsey (2025), \emph{Emergent
Introspective Awareness in Large Language Models}. Röttger et
al.~(2024), \emph{XSTest}. Turpin et al.~(2023), \emph{LMs Don't Always
Say What They Think}.

\end{document}
